\section{Materials and methods}

\subsection{Training program}
The study was conducted in a bank. The questions were prepared by the \textit{Banking Process Support
and Optimization Department}, which also decided on employee participation in the program. The
questions concerned knowledge used in everyday work at the bank. Each was a multiple-choice question
with three to five answer options, one of which was correct, while the others were
\emph{distractors}. The questions, answer options and explanations were written in Polish.

The training application started automatically on working days on employees' computers and displayed
one to four randomly selected questions per day, between 9:00 and 16:00. After each response, correct
or incorrect, it showed the correct answer with an explanation (\emph{feedback}). The same question
could return to an employee no earlier than after seven days, and the repetition schedule did not
depend on the response outcome. The program was suspended while the questions were reviewed and
updated. In the analyzed data, the median number of responses per employee on an active day was
\DailyMedian{} (IQR \DailyQOne--\DailyQThree). The same question returned after a median of
\GapMedian{} days (IQR \GapQOne--\GapQThree).

An example question presented to program participants, on cybersecurity, read:
``Which of the following situations is phishing?''. Answer options:
``1 - A DDoS attack'',
``2 - Installing antivirus software'',
``3 - Backing up data'',
``4 - A fake e-mail with a link to the bank's login page (marked as the correct answer)''.
Explanation displayed after the response:
``Phishing is a fraud method in which someone impersonates a trusted person or institution (e.g., a bank, a courier, a public office) to trick you into revealing confidential data such as passwords, card numbers or logins. Most often it looks like an e-mail from the `bank' with an urgent request to log in, a text message with a link to an `extra payment for a parcel', or a fake website deceptively similar to the real one. The key mechanism is that the fraudster creates a sense of urgency or stress so that you click without thinking. IMPORTANT: do not click suspicious links, always check the website address (URL), do not enter your login details via a link from a message, and if something looks `too urgent', it is often a scam.''.

\subsection{Data and study population}
\label{sec:dane}
The recorded employee responses were analyzed. Each record contained the question, the selected
option, the employee code, the date and time of the response, and whether the response was correct.
The analytical dataset comprised \NResponses{} responses
(Table~\ref{tab:zbior_danych}). The bank also provided a summary of employees' age and sex.

\begin{table}[H]
\centering
\caption{Characteristics of the dataset.}
\label{tab:zbior_danych}
\small
\tabelacsv{lr}{zbior_danych}
\uwagatabeli{IQR: interquartile range. A sequence is a series of attempts by one employee at one
question. The number of observed answer options may be smaller than the number of options in the
question because some options were not chosen by any employee. Excluded test responses: \NExcluded.
Excluded questions with a single observed answer option: \NExcludedItems{} (\NExcludedItemResp{}
responses). Excluded repetitions given earlier than \MinGapDays{} days after the previous attempt:
\NExcludedShortRep.}
\end{table}

Demographic data were available for \DemoMatched{} of the \DemoN{} employees (\DemoMatchedPct\%).
Women accounted for \DemoWomenPct\% of persons of known sex ($n = \DemoWomen$). The median age was
\AgeMedian{} years (IQR \AgeQOne--\AgeQThree; range \AgeMin--\AgeMax), and the largest age group was
\AgeModalGroup{} years (\AgeModalGroupPct\%; Table~\ref{tab:demografia}). Age was included in the
analysis of the first attempt and of the rate of accuracy gain; sex was not analyzed.

\begin{table}[H]
\centering
\caption{Demographic characteristics of the employees included in the analysis.}
\label{tab:demografia}
\small
\tabelacsv{lr}{demografia}
\uwagatabeli{Percentages by sex were calculated among persons of known sex, and percentages by age
group among persons of known age ($n = \DemoAgeN$). Missing data concern mainly employees absent from
the demographic summary (\DemoUnmatched). IQR: interquartile range; SD: standard deviation.}
\end{table}

In the response database, employees were identified by a random code, and the demographic summary was
linked to the recorded responses using this code. Such a code cannot be reconstructed from a list of
employees; the analyzed data were therefore fully anonymized.

\subsection{Statistical analysis}
\label{sec:analiza}
The responses of one employee to one question, ordered in time, formed a \emph{sequence}. The position
of a response in the sequence was termed the \emph{attempt number}, and the first response, that is,
the employee's first encounter with the question, the \emph{first attempt}. The \emph{interval} was the
number of days since the previous attempt in the same sequence, and \emph{accuracy} was the percentage
of correct responses. Accuracy was reported with 95\% confidence intervals (CI) obtained with the
bootstrap: employees were resampled with replacement (\BootReps{} replicates), which accounts for the
dependence between responses given by the same person. In regression models, standard errors were
clustered by employee.

\emph{First attempt.} Accuracy at the first attempt was compared across age groups. The association
with age (10-year groups) was estimated by logistic regression with question fixed effects, that is,
with a separate intercept for each question. Question difficulty was not recorded in the data, so the
fixed effects were estimated from the responses of all employees. The model thus compared employees of
different ages answering the same questions. For comparison, it was also fitted without question fixed
effects.

\emph{Changes within a sequence (RQ1).} Accuracy was calculated for successive attempt numbers;
attempts \MaxAttemptPooled{} and later were pooled. With each attempt, the number of sequences
decreased, and employees who continued responding could differ from the others, which distorts
aggregate curves \citep{murray2013}. The change was therefore also assessed in a \emph{cohort}, that
is, in sequences of at least \CohortK{} attempts, restricted to attempts 1 to \CohortK{}. Accuracy at
each of these attempts was calculated for the same sequences, and the measure of change was the
difference between attempt~\CohortK{} and attempt~1 in percentage points (pp).
The association between the rate of gain and age was tested in the cohort by logistic regression with
an interaction between $\log_2$ attempt number and age, and with employee and question fixed effects.

\emph{Outcome of the previous attempt and interval (RQ2).} Each attempt from the second onward was
assigned to one of four categories according to the outcome of that attempt and of the previous one:
a correct response after a correct one, \emph{error correction} (correct after incorrect),
\emph{forgetting} (incorrect after correct) or \emph{repeated error} (incorrect after incorrect).
Because the error-correction rate ($r$) and the forgetting rate ($f$) have different denominators, the
numbers of both events were compared, and the equilibrium level $r/(r+f)$, toward which accuracy would
tend at constant rates, was determined. The rates were compared across six interval bins:
\MinGapDays--13, 14--29, 30--59, 60--119, 120--239 and at least 240 days. The forgetting rate was also
tabulated by the number of consecutive correct responses before a given attempt and separately after an
error correction. An incorrect response at the attempt following an error correction was termed an
\emph{error relapse}, and it was checked whether the employee chose the same distractor as before the
correction. It was also checked whether the repetition schedule depended on the outcome of the
previous attempt.

\emph{Persistent errors (RQ3).} Accuracy was compared by the number of \emph{consecutive errors}, that
is, successive incorrect responses immediately preceding a given attempt, from 0 to at least \RunMax.
For two successive errors, the frequency with which the employee chose the same distractor was
calculated and compared with the percentage expected under random choice and under choice
proportional to distractor popularity in the given question. This percentage was also tabulated in
three interval bins: shorter than 30 days, 30--89 days and at least 90 days. The share of repeated
errors concentrated in the 10\% of questions with the most such errors was also calculated.

Computations were performed in Python \VerPython{} using pandas \VerPandas, NumPy \VerNumpy, SciPy
\VerScipy, statsmodels \VerStatsmodels{} and Matplotlib \VerMatplotlib.
